% ============================================================
%  Chapter 1 — What Is the Internet?
%  The Secret Life of the Internet
% ============================================================

\chapter{What Is the Internet?}

\chapteropening{%
  You tap a screen, press play, and a video appears in seconds.
  But how does it get there? Let's find out!
}
\chapterbadge{Mission Start: Meet the Internet!}

% --- Opening story ---
\section*{Meet Dot}

One afternoon, a child named Sam tapped the tablet and pressed \emph{play} on a video
about dolphins.

The video appeared instantly.

Sam looked up and wondered:

\begin{center}
  \Large\bfseries\color{InternetBlue}
  ``How did that dolphin video find me so fast?''
\end{center}

\vspace{8pt}

Out from the corner of the screen popped a tiny glowing dot with big curious eyes.

\character{Dot}{Hi! I'm Dot — a little packet of information.
  I travel through the internet every single day.
  Want to see how?  Follow me!}

\vspace{6pt}

% --- Big scene: zoom out ---
\section*{The Big Picture}

\begin{picturethis}
  Imagine standing in your bedroom holding a tablet.
  Now zoom out slowly: you see your house, then your street,
  then your whole town, then the country, then the whole planet —
  and everywhere you look there are tiny glowing lines connecting
  every home, school, library, shop, and phone.
  \textbf{That huge web of connections is the internet.}
\end{picturethis}

The internet is not one single thing you can touch.
It is a \term{network} — a system of connected devices that can
send and receive information with each other.

Right now, billions of devices around the world are connected.
Phones, tablets, laptops, smart TVs, game consoles, even some fridges —
they are all part of the same giant network.

\begin{funfact}
As of 2024, more than 5.4 billion people — over two thirds of the world's population — use the internet.
\end{funfact}

\sectionrule

% --- Simple metaphor ---
\section*{The Internet Is Like a Road System}

\begin{picturethis}
  Picture a giant system of roads.
  Every house has a driveway.
  Every driveway connects to a street.
  Every street leads to a bigger road.
  Every big road connects to a motorway.
  And motorways link whole cities and countries together.

  The internet works the same way.
  Your device connects to a small local network.
  That connects to a bigger network.
  That connects to the whole world.
\end{picturethis}

You can also think of the internet as a giant mail system.
When you send a letter, it does not fly straight to the person.
It goes to a post office, then another, then another, until it arrives.
The internet moves information the same way — in small steps, very, very fast.

\sectionrule

% --- Quick map ---
\section*{How Your Device Connects}

\begin{quickmap}
  \textbf{Step by step — from your screen to the world:}
  \begin{enumerate}[label=\textbf{\arabic*.}]
    \item Your \term{device} (tablet, phone, laptop) connects to your home network.
    \item Your home network connects to your \term{Internet Service Provider (ISP)} —
          the company that brings the internet to your home.
    \item Your ISP connects to bigger and bigger networks.
    \item Those networks connect to the whole internet.
  \end{enumerate}
\end{quickmap}

\begin{diagram}[How information starts its journey from your home to the world.]
  \begin{tikzpicture}[>=Stealth, node distance=2.5cm]
    \node[draw=InternetBlue, rounded corners, fill=SkyBlue!20, minimum width=2.6cm, minimum height=1cm] (device) {Your Device};
    \node[draw=WarmOrange!80!black, rounded corners, fill=WarmOrange!20, minimum width=2.2cm, minimum height=1cm, right=of device] (router) {Router};
    \node[draw=LeafGreen!80!black, rounded corners, fill=LeafGreen!20, minimum width=1.8cm, minimum height=1cm, right=of router] (isp) {ISP};
    \node[draw=SoftPurple!80!black, rounded corners, fill=SoftPurple!18, minimum width=3cm, minimum height=1cm, right=of isp] (cloud) {Internet Cloud};

    \draw[dashed,decorate,decoration={snake,amplitude=1.2mm,segment length=6mm},thick,InternetBlue,->] (device) -- node[below=4pt]{radio waves} (router);
    \draw[thick,InternetBlue,->] (router) -- node[below=4pt]{cable} (isp);
    \draw[thick,InternetBlue,->] (isp) -- node[below=4pt]{fibre} (cloud);
  \end{tikzpicture}
\end{diagram}

\sectionrule

% --- New words ---
\begin{newwords}
  \begin{description}[labelwidth=3cm, leftmargin=3.2cm]
    \item[\term{Internet}] A huge worldwide network that connects billions of devices.
    \item[\term{Network}]  A group of devices linked together so they can share information.
    \item[\term{Device}]   Any gadget that can connect to the internet — phone, tablet, laptop.
    \item[\term{ISP}]      Internet Service Provider — the company that connects your home to the internet.
    \item[\term{Packet}]   A tiny chunk of information travelling through the network (that's me — Dot!).
  \end{description}
\end{newwords}

\sectionrule

% --- Try It activity ---
\begin{tryit}
  \textbf{Draw your own network!}

  On a piece of paper, draw:
  \begin{itemize}
    \item Your home in the middle.
    \item Three friends' homes around it.
    \item Your school in the corner.
    \item The local library on the other side.
  \end{itemize}
  Now draw lines connecting them all together.
  Congratulations — you just drew a network!

  \emph{Can you add more buildings?  The bigger the network, the more things can talk to each other.}
\end{tryit}

\sectionrule

% --- Chapter takeaway ---
\begin{bigidea}
  The internet is a \textbf{giant network} of billions of connected devices.
  It lets computers, phones, and tablets share information with each other
  — no matter how far apart they are.
\end{bigidea}

\character{Dot}{Great start, explorer!
  Now that you know what the internet is,
  let's go and visit a website.  Follow me!}
